\section{Passive selected-port cut}
\label{sec:passive}

Consider first a finite-dimensional energy-normalized passive realization
\begin{equation}
A=-iH-\frac12K^\dagger K,
\qquad H=H^\dagger,
\label{eq:Apassive}
\end{equation}
where $K$ contains every physical radiation and loss port. Let $K_i$ and $K_o$ select two disjoint port groups with no direct cross-feedthrough. Their strictly proper transfer block is
\begin{equation}
H_{o\leftarrow i}(\nu)
=-K_o(i\nu I-A)^{-1}K_i^\dagger .
\label{eq:crossH}
\end{equation}

Define the selected-input Gramian
\begin{equation}
P_i(\tau)=\int_0^\tau
 e^{At}K_i^\dagger K_i e^{A^\dagger t}\dd t .
\end{equation}
Since $K_i^\dagger K_i\le K^\dagger K=-(A+A^\dagger)$,
\begin{equation}
\boxed{
0\le P_i(\tau)
\le I-e^{A\tau}e^{A^\dagger\tau}
\le I .
}
\label{eq:PiBound}
\end{equation}
For a stable realization, Plancherel gives
\begin{align}
\|H_{o\leftarrow i}\|_2^2
&\equiv \frac{1}{2\pi}\int_{-\infty}^{\infty}\dd\nu\,
\Tr(H_{o\leftarrow i}^\dagger H_{o\leftarrow i}) \\
&=\Tr(K_oP_iK_o^\dagger)
\le \Tr(K_o^\dagger K_o).
\end{align}
The dual output Gramian gives the complementary input bound. Therefore
\begin{equation}
\boxed{
\|H_{o\leftarrow i}\|_2^2
\le
\min\!\left[
\Tr(K_i^\dagger K_i),
\Tr(K_o^\dagger K_o)
\right].
}
\label{eq:h2cut}
\end{equation}
This is generic passive-system mathematics, not a gravitational result.

Now factor the separated gravitational link as
\begin{equation}
T(\nu)=H_B(\nu)P_g(\omega_0)H_A(\nu)
\label{eq:factor}
\end{equation}
to retained narrowband order, where $H_A$ maps selected source inputs to gravitational outputs, $H_B$ maps gravitational inputs to selected receiver outputs, and $P_g$ is the normalized free-space gravitational propagation operator. Let
\begin{equation}
\|P_g(\omega_0)\|_{\op}^2\le\eta_{\max}.
\end{equation}
Because the integrand in Eq.~\eqref{eq:gamma} is nonnegative, the band-limited spectral area is no larger than the corresponding full-line integral. Pointwise passivity gives $\|H_A(\nu)\|_{\op},\|H_B(\nu)\|_{\op}\le1$ for the selected scattering blocks. Hence
\begin{align}
\Gamma_{\rm coh}
&\le \frac{1}{2\pi}\int_{-\infty}^{\infty}\dd\nu\,
\|H_B P_g H_A\|_{\HS}^2 \\
&\le \eta_{\max}\|H_A\|_2^2
\le \eta_{\max}\Tr(K_{g,A}^\dagger K_{g,A}),
\label{eq:sourcecut}
\end{align}
where the last inequality is Eq.~\eqref{eq:h2cut} applied to the source local-to-gravitational cross block. Interchanging the roles of source and receiver gives independently
\begin{equation}
\Gamma_{\rm coh}
\le \eta_{\max}\|H_B\|_2^2
\le \eta_{\max}\Tr(K_{g,B}^\dagger K_{g,B}).
\label{eq:receivercut}
\end{equation}
Therefore
\begin{equation}
\boxed{
\Gamma_{\rm coh}
\le
\eta_{\max}
\min\!\left[
\Tr(K_{g,A}^\dagger K_{g,A}),
\Tr(K_{g,B}^\dagger K_{g,B})
\right].
}
\label{eq:gravcut}
\end{equation}

The same cut extends to a separable countably infinite modal Hilbert space when the Markov port operators are bounded and the gravitational port is Hilbert--Schmidt. If $\mathcal T(t)$ is the passive contraction semigroup,
\begin{equation}
P_i(\tau)=\int_0^\tau
\mathcal T(t)K_i^\dagger K_i\mathcal T^\dagger(t)\dd t
\end{equation}
obeys the operator inequality $0\le P_i(\tau)\le I$. The monotone strong limit therefore satisfies $0\le P_i\le I$. Section~\ref{sec:resource} shows that the retained gravitational matter resource itself implies $\Tr(K_g^\dagger K_g)<\infty$, making $K_g$ Hilbert--Schmidt and the operator-valued $H_2$ norm finite. This covers countably infinite \emph{bounded-port} modal sectors; arbitrary unbounded boundary-control PDE ports require separate admissibility analysis \cite{BarasBrockett1975,Opmeer2013}.
